%% Chapter 6 -----------------------------------------------------------------
\chapter{The Sidebar Panels}
\label{ch:sidebar}
\index{sidebar}

The sidebar on the left of the main window groups the controls used most often while
working on a spectrum. It contains eight sections, from top to bottom:
\gui{Timeline}, \gui{Background Subtraction}, \gui{Noise Clipping Thresholds}, \gui{Units},
\gui{Axis}, \gui{Graph Properties}, \gui{Analysis Summary} and \gui{Ruler Measurement}.
Each section is a card with a clickable header that collapses or expands it; the state of
every card is remembered between sessions. The whole sidebar can be hidden with the arrow
handle on the left edge of the plot.

\section{Timeline}
\label{sec:timeline}
\index{timeline}

A loaded dataset knows which station, focus codes and observation times it came from. The
\gui{Timeline} panel (Figure~\ref{fig:timeline}) lists them and extends or trims the
dataset in place, without returning to the downloader.

\screenshot[0.45\linewidth]{main-window/timeline}{The \gui{Timeline} panel.}{fig:timeline}{Timeline panel with two observations.}

\begin{description}
  \item[◀ Prev / Next ▶] Fetch the observation before the first, or after the last, and
    time-combine it into the spectrum (\kbd{Ctrl+Shift+Left} and \kbd{Ctrl+Shift+Right},
    also in the \menu{Edit} menu).
  \item[Step selector (\(\times\)1, \(\times\)2, \(\times\)4, \(\times\)8)] The number of
    observations added per click.
  \item[Observation list] One row per observation, giving its start time and focus code(s).
  \item[Trim Start / Trim End] Remove the earliest or the latest observation. Trimming down
    to a single observation returns to a plain single-file load.
  \item[Status line] The station, the covered time range and the number of observations,
    for example \gui{BIR · 08:00–08:15 · 2 observations}. When a direction is unavailable
    the panel says why, for example the end of the archive day or a timestamp that supplies
    only part of the focus-code set.
\end{description}

Adjacent observations are looked for first among the loaded file's neighbors on disk, and
only then in the e-CALLISTO archive, so a set that has already been downloaded extends
without network access. Day listings are cached during the session. With
\menu{Edit > Prefetch Adjacent Observations}\index{timeline!prefetch} ticked (the default), neighboring files are
downloaded in the background so that the next step is instant; untick it to avoid
unrequested archive traffic.

Extending and trimming preserve your work:
\begin{itemize}
  \item annotations, the ruler measurement and drift picks are kept, and their time
    coordinates shift with the axis when an earlier observation is prepended;
  \item the active background subtraction and RFI cleaning are recalculated over the
    longer array rather than returning the view to the raw data;
  \item the display thresholds are kept;
  \item every extension or trim is a single undo\index{undo} step (\kbd{Ctrl+Z}), restoring the arrays,
    the source list and the annotations together.
\end{itemize}
After an extension or trim the plot is rescaled to show the full spectrum.

\section{Background Subtraction}
\label{sec:background}
\index{background subtraction}

The receiver background varies strongly from channel to channel, producing horizontal
stripes in a raw spectrum. The \gui{Background Subtraction} panel
(Figure~\ref{fig:bg-sub}) removes it.

\screenshot[0.45\linewidth]{main-window/bg_sub}{The \gui{Background Subtraction} panel.}{fig:bg-sub}{Background Subtraction panel.}

Choose a \gui{Method} and click \gui{Subtract Background}\index{background subtraction!Subtract Background}. The background is computed for
each frequency channel over the whole observation time:
\begin{description}
  \item[Mean] subtracts the mean of each frequency channel; the result stays in the raw
    intensity unit (Digits).
  \item[Median] subtracts the median of each frequency channel; the result stays in
    Digits. The median is less affected by bursts and RFI than the mean.
  \item[Median (dB)] subtracts the median of each channel and converts the result to
    decibels above that background, following the CALLISTO Plotutil\index{Plotutil} recipe used by the
    batch processor. This method is recommended for analysis.
\end{description}

The status line below the button names the method in use, for example
\gui{Applied: Median (dB)}\index{background subtraction!Median (dB)}. The subtraction is always computed from the raw data, so
choosing another method and clicking again replaces the previous result rather than
stacking on it. Re-applying the subtraction also clears any RFI cleaning or isolated burst
that was derived from the previous result. While a Median (dB) result is shown, the
intensity unit is locked to dB because the data are already in decibels.
\menu{Edit > Reset to Raw} removes the subtraction.

\section{Noise Clipping Thresholds}
\label{sec:thresholds}
\index{noise clipping thresholds}\index{display thresholds}

The \gui{Noise Clipping Thresholds} panel (Figure~\ref{fig:noise-clipping}) sets the
limits of the color scale.

\screenshot[0.45\linewidth]{main-window/noise_clipping}{The \gui{Noise Clipping Thresholds} panel.}{fig:noise-clipping}{Noise Clipping Thresholds panel.}

\begin{description}
  \item[Lower Threshold / Upper Threshold] High-resolution sliders for the lower and upper
    limits of the color scale (Vmin and Vmax). The display updates live while you drag;
    no Apply button is needed. The value card beside each label shows the limit. When the
    intensity is shown in Digits, a second line gives the corresponding dB value.
  \item[Logarithmic Threshold Scale] Makes the sliders logarithmic for finer control near
    zero.
\end{description}

The thresholds change the contrast only; the data are never clipped or modified. With
both sliders at zero the color scale spans the whole data range. Limits chosen on the raw
data are reset when you subtract the background, because they would saturate the new
scale. \menu{Processing > Presets > Raw FITS Percentile (5-98\%)}\index{presets!Raw FITS Percentile} sets the limits from the
distribution of the data on show (Section~\ref{sec:presets}).

\section{Units}
\label{sec:units}
\index{units}

The \gui{Units} panel (Figure~\ref{fig:units}) selects the units of the intensity and time
axes.

\screenshot[0.45\linewidth]{main-window/units}{The \gui{Units} panel.}{fig:units}{Units panel.}

\begin{description}
  \item[Intensity: Digits / dB] \gui{Digits} shows raw CALLISTO ADC values as stored in
    the file. \gui{dB} converts them with the CALLISTO receiver scale (see the glossary
    entry \emph{Digit}).
  \item[Time: Seconds / UT] \gui{Seconds} measures time from the start of the loaded data;
    \gui{UT} shows Universal Time. UT requires the observation start time
    (\texttt{TIME-OBS}\index{TIME-OBS keyword}) in the FITS header.
\end{description}
Switching units does not modify or lose data; it only changes the axis labels, the
color-bar label and the cursor readout.

\section{Axis}
\label{sec:axis}
\index{frequency axis!logarithmic}

The \gui{Axis} panel (Figure~\ref{fig:axis}) switches the frequency axis between
\gui{Linear}\index{height--time fit!order} (proportional to MHz) and \gui{Log} (proportional to \(\log_{10}\) MHz).

\screenshot[0.45\linewidth]{main-window/axis}{The \gui{Axis} panel.}{fig:axis}{Axis panel.}

A logarithmic axis spreads out the decametric end of the band, where Type~II and
Type~III bursts spend most of their drift, and places a fundamental--harmonic pair at a
constant separation. Ticks are anchored to the decades and labeled in MHz (for example 20,
30, 40, 50, 60, 70, 80 across a 20--80~MHz band), thinning as the span widens. Both
renderers support the logarithmic axis, and annotations, the ruler, light-curve picks and
drift points behave identically in both scales because their coordinates remain in MHz.
The choice is remembered between sessions. The STEREO/SWAVES panel always uses its own
logarithmic axis (Section~\ref{sec:swaves}).

\section{Graph Properties}
\label{sec:graph-properties}
\index{graph properties}\index{color map}

The \gui{Graph Properties} panel (Figure~\ref{fig:graph-properties}) controls the
appearance of the plot. It becomes active once a file is loaded. The settings apply to
the on-screen plot and to exported figures and reports.

\screenshot[0.45\linewidth]{main-window/graph_properties}{The \gui{Graph Properties} panel.}{fig:graph-properties}{Graph Properties panel.}

\begin{table}[!htbp]
  \centering
  \caption{Controls of the \gui{Graph Properties} panel.}
  \label{tab:graph-properties}
  \small
  \begin{tabularx}{\linewidth}{@{}P{0.16\linewidth}P{0.24\linewidth}L@{}}
    \toprule
    \thead{Group} & \thead{Control} & \thead{Function} \\
    \midrule
    Appearance & Colormap & \gui{Custom} (blue--red--yellow), \gui{viridis}, \gui{plasma}, \gui{inferno}, \gui{magma}, \gui{cividis}, \gui{turbo}, \gui{RdYlBu}, \gui{jet}, \gui{cubehelix} or \gui{bone\_r}. \\
               & Font family\index{graph properties!font} & Typeface of all plot text; \gui{Default} or any installed font. \\
    Text & Graph title & A custom title; leave empty for the default title. \\
         & Remove Titles & Hide the plot title. \\
    Font sizes & Tick labels (px) & 6--60~px (default 11). \\
               & Axis labels (px) & 6--60~px (default 12). \\
               & Title (px) & 6--80~px (default 14). \\
    Text style & Title, Axis labels, Tick labels & \gui{Bold} and \gui{Italic} for each text element. \\
    \bottomrule
  \end{tabularx}
\end{table}

\section{Analysis Summary}
\label{sec:analysis-summary}
\index{analysis summary}

The \gui{Analysis Summary} panel (Figure~\ref{fig:analysis-summary}) is a read-only summary
of the current analysis session: the fitted equation, \(R^2\), the fold number\index{fold number} and the
average shock speed and height from the Burst Analyzer (Section~\ref{sec:analyzer}). It
also reports when a saved session has been restored. Before any analysis it reads
\gui{No analysis session loaded.}

\screenshot[0.45\linewidth]{main-window/analysis_summary}{The \gui{Analysis Summary} panel.}{fig:analysis-summary}{Analysis Summary panel.}

\section{Ruler Measurement}
\label{sec:ruler-panel}
\index{ruler}

The \gui{Ruler Measurement} panel (Figure~\ref{fig:ruler-panel}) shows the result of the
ruler tool, started with \menu{Analysis > Ruler Measurement}
(Section~\ref{sec:ruler}): the two points \gui{P1} and \gui{P2} (time and frequency), the
\gui{Duration}, the \gui{Frequency drift} and the \gui{Slope} in MHz/s. \gui{Copy} copies
the readout to the clipboard and \gui{Clear} removes the measurement.

\screenshot[0.45\linewidth]{main-window/ruler_measurements}{The \gui{Ruler Measurement} panel.}{fig:ruler-panel}{Ruler Measurement panel.}
